\chapter{Conclusion}
\label{chap:conclusion}

This thesis proposed an end-to-end Personalized Travel Guide Assistant for
Vietnam based on Retrieval-Augmented Generation (RAG). The system was designed
to address three connected challenges: understanding multi-turn travel
requests, retrieving relevant and sufficiently complete tourism evidence, and
generating responses that are both grounded in the available sources and
adapted to the individual user.

The proposed architecture integrates an offline tourism knowledge-construction
workflow with an inspectable online query-processing pipeline. Tourism webpages
and documents are cleaned, divided using structure-aware chunking, enriched
with controlled metadata, transformed into dense embeddings, and stored in a
PostgreSQL knowledge base with vector-search support. At query time, the system
performs request routing, conversational query rewriting, structured query
parsing, retrieval planning, hybrid retrieval, CrossEncoder reranking,
evidence-coverage checking, bounded recovery retrieval, and readiness-gated
answer generation.

A central contribution of this work is the separation of different types of
context and query information. Persistent user memory stores durable
preferences, such as preferred activities, travel styles, budget level, answer
length, and conditions to avoid. Temporary conversation state stores
trip-specific information, including the current destination, dates, duration,
selected places, and temporary constraints. Similarly, the query-understanding
stage distinguishes intent from operation and separates explicit user
constraints from semantic retrieval facets. These distinctions reduce the risk
of treating temporary requirements as permanent preferences or applying soft
personalization signals as restrictive retrieval filters.

The retrieval pipeline combines dense semantic retrieval and BM25 lexical
retrieval using Reciprocal Rank Fusion. Small metadata, preference, geographic,
and freshness signals refine the ranking, while a CrossEncoder reranks the
strongest fused candidates. For complex requests, a bounded planner decomposes
the query into focused evidence requirements and retrieval tasks. The evidence
checker then determines whether each requirement is covered, partially covered,
or not covered. When evidence remains incomplete, the system performs up to
three focused recovery searches and evaluates the evidence again.

The proposed design also distinguishes retrieval failure from corpus
limitations. Missing-evidence diagnostics identify whether a requirement was
not probed, whether the corpus likely lacks the requested information, whether
potentially useful chunks were lost during reranking, or whether selected
evidence remains insufficient. Based on these results, the system assigns a
complete, partial, or insufficient answer-readiness mode. This prevents the
answer model from silently replacing missing evidence with unsupported internal
knowledge and enables partial answers to state their limitations explicitly.

To evaluate the complete pipeline, this thesis defines a multi-stage evaluation
protocol covering query understanding, retrieval, and final-answer quality.
The dataset contains five synthetic users with different travel preferences,
two multi-turn conversations per user, and ten query turns per conversation,
for a total of 100 evaluated queries. Query understanding is evaluated using
intent accuracy, operation accuracy, query-constraint F1, and retrieval-facet
F1. Retrieval quality is evaluated using graded chunk relevance, nDCG@5, and
Precision@5. Final responses are evaluated for correctness, faithfulness,
personalization adherence, and completeness.

An independent LLM is used to provide scalable initial judgments, while human
review is retained for reference-label validation and confirmation of
questionable or representative cases. The experimental trace also records
latency, token usage, estimated cost, model identifiers, prompt hashes,
pipeline version, and corpus fingerprint. Consequently, evaluation results can
be connected to the precise system and knowledge-base configuration under
which they were produced.

The experimental findings presented in this thesis show that final-answer
quality cannot be explained by a single component score. A relevant top-ranked
chunk does not guarantee that every part of a complex travel request is
covered, while a lower final-answer score may originate from missing corpus
information, query-parsing error, insufficient task decomposition, ranking
failure, incorrect coverage assessment, or generation failure. Stage-level
evaluation and pipeline traces are therefore necessary for interpreting the
behavior of a personalized RAG system.

% Replace the following commented paragraph after the official 100-query run
% with the final aggregate results reported in the evaluation chapter.
% The official experiment achieved an Intent Accuracy of [VALUE], an Operation
% Accuracy of [VALUE], a Query Constraint F1 of [VALUE], and a Retrieval Facet
% F1 of [VALUE]. Retrieval achieved an nDCG@5 of [VALUE] and Precision@5 of
% [VALUE]. The mean final-answer scores were [VALUE] for correctness, [VALUE]
% for faithfulness, [VALUE] for personalization adherence, and [VALUE] for
% completeness. These results indicate [FINAL EVIDENCE-BASED INTERPRETATION].

\section*{Limitations}

Several limitations remain. First, system performance depends on the scope,
quality, freshness, and metadata consistency of the tourism corpus. Information
such as current prices, schedules, temporary closures, dietary safety, and
accommodation details may be absent or become outdated. Recovery retrieval can
find overlooked evidence, but it cannot create information that does not exist
in the knowledge base.

Second, complex travel requests may contain more aspects than can be represented
by the bounded retrieval plan and final evidence set. Although destination
comparison receives balancing, general multi-aspect requests can still be
dominated by evidence for one task or destination. Global CrossEncoder
reranking does not always guarantee balanced coverage across every planned
requirement.

Third, LLM-based query parsing, planning, coverage checking, response
generation, and evaluation remain probabilistic. Controlled vocabularies and
deterministic normalization reduce malformed or inconsistent outputs but cannot
eliminate them completely. Synthetic reference labels can also be ambiguous,
particularly for overlapping intents and semantic facets. Independent judging
reduces circularity but does not remove judge bias or variance, making human
validation necessary.

Finally, the current evaluation is performed on a bounded synthetic multi-turn
dataset. The dataset supports controlled comparison and repeatability, but it
does not fully represent the linguistic diversity, incomplete information,
unexpected behavior, and changing preferences of real users. Conclusions must
therefore be interpreted within the evaluated corpus, user profiles, prompts,
models, and pipeline version.

\section*{Future Work}

Future directions of this research include:

\begin{enumerate}[label=(\roman*)]
    \item \textbf{Expanding and refreshing the tourism corpus.} Future work
    should add more official and destination-specific sources, improve coverage
    of accommodation, transportation, food safety, accessibility, prices, and
    opening information, and introduce scheduled freshness validation.

    \item \textbf{Improving requirement-aware retrieval.} Retrieval and
    reranking could allocate evidence quotas per planned requirement or apply
    task-aware diversification so that one high-scoring aspect does not dominate
    a complex itinerary, comparison, or recommendation request.

    \item \textbf{Strengthening retrieval planning.} Planner evaluation could
    be extended to measure under-decomposition and over-decomposition directly.
    Adaptive task limits could then be selected according to query complexity
    rather than using one fixed maximum for every request.

    \item \textbf{Introducing live-data tools.} Time-sensitive information,
    such as weather, transport schedules, ticket prices, and temporary closures,
    could be obtained from verified external services and clearly separated
    from the static RAG corpus.

    \item \textbf{Extending personalization and memory controls.} Future
    versions could allow users to inspect, correct, deactivate, or delete stored
    preferences and could evaluate memory precision, temporal decay, conflicts,
    and privacy more extensively.

    \item \textbf{Conducting larger human and real-user evaluations.} The
    synthetic benchmark should be supplemented with independently authored
    queries, multiple human relevance assessors, inter-annotator agreement, and
    usability studies involving real travel-planning sessions.

    \item \textbf{Optimizing latency and cost.} Stage-level telemetry can be
    used to study model routing, caching, batching, smaller specialized models,
    early stopping, and conditional recovery so that quality improvements are
    balanced against response time and operational cost.

    \item \textbf{Supporting multilingual interaction.} Vietnamese and other
    languages could be evaluated explicitly, including cross-lingual retrieval,
    entity normalization, and preservation of user constraints during query
    rewriting and translation.
\end{enumerate}

Overall, this work demonstrates a practical methodology for building a
personalized and evidence-aware travel assistant in which query understanding,
memory, retrieval, coverage, and answer generation can be examined separately.
The proposed system provides a reproducible foundation for future research on
grounded travel recommendation, multi-turn personalization, adaptive retrieval,
and trustworthy knowledge-gap handling.
